\scanpage{36}
\mspage{28}
% L1–6
de $\Delta$ \deleted{est} s'identifie au gpe des autom de $(T\longrightarrow I)$, une fois choisi un iso. entre les deux fibres de $T\to I$, il s'identifie au produit $\tfrac12$ direct
\[
\underset{\substack{\shortparallel\\\mathfrak S_2}}{\unsure{\mathfrak S|\mathfrak G}_{\unsure{I|2}}}\cdot\prod_{i\in I}\mathfrak S_{T_i}
\]
de cardinal 72 (\uline{égal} au nb de racines\ldots).

% L7–15
Si on a un plongt $\Delta_0\xrightarrow{i} C$, \deleted{les} \inserted{le gpe des} autom de $(C,i(\Delta_0)=\Delta)$ s'applique sur celui de $\Delta$, le noyau étant d'ordre 6, isomorphe : $\mathfrak S_3$. On va associer à $i(\Delta)$ un ens $V^*$, \unsure{très strict|trois élémts|triple}, qui sera canoniqut isom, pour tt $u\in i(\Delta)$, à \deleted{\unsure{un}} l'ens \inserted{$V^*(u)$} des trois triangles de $C$ de sommet $u$, distinct des 2 triangles $\in\Delta$. Il faut donc établir

% L16–20; left rule along the four lines of the lemma.
\uline{Lemme} \unsure{1}\quad $\exists!$ syst. transitif d'iso. entre les $V^*(u)$ $(u\in\Delta)$ tels que si $x\in C-\Delta$, et si $x$ est lié à $u,v\in\Delta$, alors les triangles $\operatorname{tr}(u,x)\in V^*(u)$ et $\operatorname{tr}(v,x)\in V^{\unsure{1|*}}(v)$ se correspondent

% L21–25
Utiliser l'épinglage des p. 8 et 9

On trouve que
\[
\operatorname{Aut}(C,\Delta)\longrightarrow
\underset{\substack{\shortparallel\\\mathfrak S_2\cdot(\mathfrak S_3\times\mathfrak S_3)}}{\operatorname{Aut}(\Delta)}
\times\underset{\substack{\shortparallel\\\mathfrak S_3}}{\operatorname{Aut}(\unsure{W|V}^*)}
\]
est un iso. Il faut donc expliciter comment on reconstitue $(C,\Delta)$
% archive sheet number: 35
